2222\hspace{3mm}

\centerline{\bf \Huge Домашнее задание.}
\centerline{\bf \Huge Теорема Без У.}

\begin{flushright}
    \scriptsize{Неважно, возможно это или нет. \\Я это сделаю, потому что хочу сделать.\\ Монки Д. Луффи}
\end{flushright}

\problem{\textbf{1!}} При каком значении $a$ многочлен $P(x) = x^{1991} + ax^{1922} + 69$ делится на $x-1$?

\problem{\textbf{2!}} Разделите многочлен $x^5 + 5x^3 + 6$ на многочлен $x^2 + 2x + 3$, найдите частное и остаток.

\textit{Примечание.} Нужно привести полные вычисления, а не просто ответ.


\problem{3} a) Докажите, что если число $a$ является корнем многочлена $P(x)$, то $P(x) =\\ (x-a)Q(x)$, причем степень многочлена $Q(x)$ на 1 один меньше, чем у $P(x)$.

\noindentб) Докажите, что у многочлена степени $n \geq 0$ не более чем $n$ корней.

\problem{4} Олег попросил Эльзяту назвать любой многочлен $P(x)$ с \textbf{натуральными} коэффициентами, а затем назвал такое число $k$, что $P(k), P(k+1), \dotsc, P(k+2024)$ являются составными числами. Эльзята подумала, что eё брату просто повезло и придумала ещё один многочлен. Но Олег и тут смог найти подходящее число $k$. Правда ли, что Олег всегда сможет находить такое число, какой бы многочлен не загадывала Эльзята или ему просто дважды повезло?